\chapter{The Discrete Logarithm Problem and ElGamal}

In the chapter on Diffie-Hellman you met a strange one-way street. Given a generator $g$ of some
group and an exponent $x$, computing $g^x$ is cheap -- a few hundred multiplications even for
numbers thousands of bits long. Going backwards, recovering $x$ from $g^x$, appears to cost an
astronomical amount of work. That backwards problem has a name, the \emph{discrete logarithm
problem}, and almost everything in public-key cryptography that is not built on factoring is built
on it instead: Diffie-Hellman key exchange, the ElGamal encryption scheme we develop here, the DSA
and Schnorr signatures of later chapters, and -- with the group swapped out for points on an
elliptic curve -- the entire ECC ecosystem that protects your phone calls and web traffic today.

This chapter does three things. First, we pin down exactly what the discrete log problem is, in a
general group and in the concrete group $\Zns{p}$, and we trace tiny numerical examples by hand so
the asymmetry between forward and backward is something you have felt, not just read. Second, we
study the two classic algorithms that attack it -- baby-step giant-step and Pollard's rho -- because
``hard'' in cryptography never means ``impossible,'' and the precise cost of the best attack is what
sets your key sizes. Third, we build ElGamal encryption on top of the hardness assumption, prove it
secure under the Decisional Diffie-Hellman assumption, discover its surprising homomorphic property,
and work out how to choose groups that do not quietly hand the attacker an easy win.

\chapterline

\section{The discrete logarithm, precisely}
\label{ch20:dlp}

Fix a finite cyclic group $\G$ of order $n$, written multiplicatively, and a generator $g$. That $g$
generates $\G$ means every element of $\G$ equals $g^k$ for some integer $k$, and the map
\[
  k \;\longmapsto\; g^k, \qquad \Zn{n} \longrightarrow \G,
\]
is a bijection (a group isomorphism, in fact, since $g^{a+b}=g^a g^b$). The \emph{discrete
logarithm} of an element $h \in \G$ to base $g$ is the exponent that produces it.

\begin{definition}[Discrete logarithm]
\label{ch20:def-dlog}
Let $\G = \langle g \rangle$ be cyclic of order $n$. For $h \in \G$, the discrete logarithm
$\dlog_g h$ is the unique residue $x \in \Zn{n}$ with $g^x = h$. The \emph{discrete logarithm
problem} (DLP) is: given $\G$, $g$, and $h$, compute $x$.
\end{definition}

The word ``logarithm'' is exactly the right analogy. Over the real numbers, $\log_b y$ is the
exponent you raise $b$ to in order to get $y$, and it is easy because the real exponential is smooth
and monotone -- you can binary-search it. The word ``discrete'' is the warning. In a finite group
the powers $g^0, g^1, g^2, \dots$ do not march steadily upward; they hop around the group with no
useful order. There is nothing to binary-search. Knowing $g^{x}$ is just under $h$ tells you nothing
about whether $x$ is close to $\dlog_g h$, because ``just under'' is meaningless once you reduce mod
$p$.

\begin{intuitionbox}
Think of a clock with $p-1$ positions and a hand that, each tick, jumps forward by a multiplicative
step. Telling someone the current position is easy if you know how many ticks have passed. But if I
only show you where the hand is pointing, and the step size scrambles the positions, you have no
shortcut: you seem forced to try tick counts one by one until the hand lands where I said. Forward is
a walk; backward is a search.
\end{intuitionbox}

\subsection{The concrete case: $\Zns{p}$}

The workhorse group for classical discrete-log cryptography is $\Zns{p}$, the nonzero residues mod a
prime $p$ under multiplication. This group is cyclic of order $p-1$ (a fact proved in the finite
fields chapter; every finite subgroup of the multiplicative group of a field is cyclic), so it has
generators. A generator of $\Zns{p}$ is traditionally called a \emph{primitive root} mod $p$.

\begin{example}[A primitive root mod 7]
\label{ch20:ex-modp}
Take $p = 7$, so $\Zns{7} = \{1,2,3,4,5,6\}$ has order $6$. Try $g = 3$ and list its powers:
\[
  3^0 = 1,\quad 3^1 = 3,\quad 3^2 = 2,\quad 3^3 = 6,\quad 3^4 = 4,\quad 3^5 = 5 \pmod 7,
\]
and $3^6 = 1$, closing the loop. The powers $1,3,2,6,4,5$ run through every nonzero residue, so $3$
is a primitive root. Now read the table backwards. We have $3^4 = 4$, so $\dlog_3 4 = 4$. And
$3^5 = 5$, so $\dlog_3 5 = 5$. There is no pattern relating the value $4$ to its log $4$, or $5$ to
its log $5$; the coincidence is an artifact of this tiny example. The sequence of logs $0,1,2,3,4,5$
corresponds to values $1,3,2,6,4,5$ -- thoroughly shuffled.
\end{example}

Even at $p=7$ you can see the difficulty in miniature. If I tell you $g=3$ and $h=6$ and ask for
$\dlog_3 6$, you cannot reason ``$6$ is close to $7$ so the exponent is large.'' You have to walk the
table: $3,2,6$ -- found it at exponent $3$. With $p$ a 2048-bit prime the table has about $2^{2048}$
rows and you cannot walk it in the lifetime of the universe.

\begin{focusbox}
\textbf{Two flavors of ``hard.''} The DLP is the \emph{search} problem: recover the exponent.
Cryptography often needs a weaker-looking but sufficient assumption, the \emph{Computational
Diffie-Hellman} (CDH) problem: given $g$, $g^a$, $g^b$, compute $g^{ab}$ -- without necessarily
finding $a$ or $b$. And ElGamal's proof will rest on a still weaker assumption, \emph{Decisional
Diffie-Hellman} (DDH): merely \emph{distinguish} $g^{ab}$ from a random group element. These three
form a hierarchy, $\text{DDH} \le \text{CDH} \le \text{DLP}$: an algorithm for the harder one solves
the easier ones. We will use each where it fits.
\end{focusbox}

\subsection{Why we believe it is hard}

There is no proof that the discrete log problem is hard -- such a proof would resolve $\mathrm{P}$
versus $\mathrm{NP}$-style questions far beyond present mathematics. What we have is the experience
of decades. The DLP in $\Zns{p}$ has been studied since the 1970s, attacked by the same number-field
sieve machinery that attacks factoring, and the best known algorithms remain \emph{subexponential}
but not polynomial: roughly
\[
  L_p[1/3, c] \;=\; \exp\!\Big( (c+o(1))\,(\ln p)^{1/3} (\ln\ln p)^{2/3} \Big)
\]
operations, the same shape as the cost of factoring an integer of similar size. This is why DLP-mod-
$p$ keys and RSA keys are comparable in length: 2048-bit and 3072-bit primes are the current
mainstream. On elliptic-curve groups, by contrast, no subexponential attack is known, the cost stays
fully exponential, and that is exactly why a 256-bit elliptic-curve key matches a 3072-bit
mod-$p$ key in strength. We will see the elliptic-curve story in the next chapter; here we develop
everything over $\Zns{p}$ where the arithmetic is concrete.

\begin{historybox}
The hardness of discrete logs is the secret sauce of the 1976 Diffie-Hellman paper, \emph{New
Directions in Cryptography} -- the paper that opened public-key cryptography. Whitfield Diffie and
Martin Hellman did not have a public-key \emph{encryption} scheme; they had a key-\emph{exchange}
protocol and a vision. It was Taher ElGamal, in his 1985 paper ``A Public Key Cryptosystem and a
Signature Scheme Based on Discrete Logarithms,'' who turned the Diffie-Hellman idea into a full
encryption scheme and a signature scheme. ElGamal later became chief scientist at Netscape and had a
direct hand in SSL, the ancestor of TLS, so the line from this chapter to the padlock in your browser
is unusually short.
\end{historybox}

\section{Attacking the DLP: generic algorithms}
\label{ch20:generic}

Before trusting a group we should know the best attacks on it. Two algorithms work in \emph{any}
cyclic group, using only the group operation -- they do not exploit the special structure of
$\Zns{p}$. These are the \emph{generic} attacks, and they set a floor that no DLP-based scheme can
beat. A theorem of Shoup says that a purely generic algorithm needs about $\sqrt{n}$ group
operations in a group of order $n$, and both algorithms below hit that bound.

The headline consequence: \textbf{the security of a DLP group is governed by the largest prime factor
of its order, and the cost of breaking it is about the square root of that prime.} If $n$ has a
small prime factor, the problem partially collapses (we will see why via Pohlig-Hellman in the
exercises). So we want $n$ to have a large prime factor -- ideally $n = 2q$ with $q$ a large prime,
the ``safe prime'' setup.

\subsection{Baby-step giant-step}

Daniel Shanks's baby-step giant-step (BSGS) is a beautiful time-memory trade-off. We want $x$ with
$g^x = h$, where $0 \le x < n$. Write $m = \ceil{\sqrt{n}}$ and split $x$ in base $m$:
\[
  x \;=\; i \cdot m + j, \qquad 0 \le i < m,\quad 0 \le j < m.
\]
Such a decomposition always exists because $i,j$ together range over $0 \dots m^2 - 1 \ge n-1$. Now
rearrange the equation $g^{im+j} = h$:
\[
  g^{j} \;=\; h \cdot (g^{-m})^{i}.
\]
The left side depends only on $j$ (the ``baby steps,'' small increments of $g$); the right side only
on $i$ (the ``giant steps,'' big leaps of size $m$). The algorithm:

\begin{enumerate}
  \item \textbf{Baby steps.} Compute $g^0, g^1, \dots, g^{m-1}$ and store each in a hash table keyed
    by the group element, with the exponent $j$ as the value. Cost: $m$ multiplications, $m$ table
    entries.
  \item \textbf{Giant steps.} Compute $u = g^{-m}$ once. Then for $i = 0, 1, 2, \dots$ form
    $h \cdot u^{i}$ and look it up in the table. The moment $h \cdot u^{i}$ equals some stored
    $g^{j}$, you have $g^{j} = h \cdot g^{-im}$, hence $x = im + j$.
\end{enumerate}

Total work is $O(\sqrt{n})$ group operations and $O(\sqrt{n})$ storage. That square root is huge:
for $n \approx 2^{256}$ you would need $2^{128}$ operations and $2^{128}$ words of memory, both
hopelessly out of reach. But it is also a dramatic improvement over the $2^{256}$ of brute force,
and it is why the relevant prime needs to be twice the bit-length of your security target.

\begin{example}[BSGS by hand, $p=29$]
\label{ch20:ex-bsgs}
Let $p = 29$, so $n = 28$, and take the primitive root $g = 2$. We want $\dlog_2 17$. Set
$m = \ceil{\sqrt{28}} = 6$.

\emph{Baby steps} -- powers $2^j \bmod 29$ for $j=0\dots5$:
\[
  2^0=1,\ \ 2^1=2,\ \ 2^2=4,\ \ 2^3=8,\ \ 2^4=16,\ \ 2^5=3.
\]
Store the table $\{1{:}0,\ 2{:}1,\ 4{:}2,\ 8{:}3,\ 16{:}4,\ 3{:}5\}$.

\emph{Giant step size.} We need $g^{-m} = 2^{-6} \bmod 29$. First $2^{6} = 64 \equiv 6$, and
$6^{-1} \bmod 29 = 5$ (since $6 \cdot 5 = 30 \equiv 1$). So $u = 2^{-6} \equiv 5$.

\emph{Giant steps} -- compute $17 \cdot 5^{i} \bmod 29$ and check the table:
\[
\begin{aligned}
  i=0:\ & 17 \cdot 1 = 17 \quad(\text{not in table})\\
  i=1:\ & 17 \cdot 5 = 85 \equiv 27 \quad(\text{no})\\
  i=2:\ & 27 \cdot 5 = 135 \equiv 19 \quad(\text{no})\\
  i=3:\ & 19 \cdot 5 = 95 \equiv 8 \quad(\textbf{hit! } 8 \mapsto j=3).
\end{aligned}
\]
So $i = 3$, $j = 3$, and $x = im + j = 3\cdot 6 + 3 = 21$. Check: $2^{21} \bmod 29$. We have
$2^{5}=3$, $2^{10}=9$, $2^{20}=81\equiv 81-58=23$, and $2^{21}=46\equiv 17$. Correct,
$\dlog_2 17 = 21$.
\end{example}

\begin{connectionbox}
BSGS is the same idea as the meet-in-the-middle attack on double encryption from the block-cipher
chapter, and the same idea as a two-sided search in the algorithms companion. You trade an
exponential-time one-sided search for two square-root-sized searches that ``meet'' through a hash
table. The cost analysis -- $O(\sqrt{n})$ time, $O(\sqrt{n})$ space, balanced when both halves are
$\sqrt{n}$ -- is a textbook time-memory trade-off; see the discussion of meet-in-the-middle in the
algorithms book.
\end{connectionbox}

\subsection{Pollard's rho: square-root time, constant memory}

BSGS has a weakness: it needs $\sqrt{n}$ \emph{memory}, and at cryptographic sizes memory is the
binding constraint long before time is. Pollard's rho method achieves the same $O(\sqrt{n})$
\emph{time} with only a constant amount of memory, by replacing the stored table with a cleverly
randomized walk and the birthday paradox.

The idea: define a deterministic but ``random-looking'' function $f$ on the group that, at each step,
multiplies the current element by either $g$, by $h$, or squares it, while carefully tracking the
exponents of $g$ and $h$ that produced the current element. So at every point of the walk we know
integers $a,b$ with
\[
  x_k \;=\; g^{a_k} h^{b_k}.
\]
The walk wanders through the group. Because the group is finite, the sequence $x_0, x_1, x_2, \dots$
must eventually repeat, and once it repeats it cycles forever -- the trajectory looks like the Greek
letter $\rho$, a tail leading into a loop (Figure~\ref{ch20:fig-rho}), which is where the method gets
its name. By the birthday
bound (chapter on hashing), a collision $x_i = x_j$ with $i \ne j$ shows up after only about
$\sqrt{\pi n / 2}$ steps. A collision means
\[
  g^{a_i} h^{b_i} = g^{a_j} h^{b_j}
  \;\Longrightarrow\;
  g^{a_i - a_j} = h^{b_j - b_i}
  \;\Longrightarrow\;
  g^{a_i - a_j} = g^{x (b_j - b_i)},
\]
so, reading exponents modulo $n$,
\[
  a_i - a_j \equiv x\,(b_j - b_i) \pmod n,
  \qquad\text{hence}\qquad
  x \equiv (a_i - a_j)\,(b_j - b_i)^{-1} \pmod n,
\]
provided $b_j - b_i$ is invertible mod $n$ (true whenever $n$ is prime and $b_j \ne b_i$). We solved
for the discrete log without ever building a table.

To detect the collision without storing the whole walk, Pollard borrows Floyd's tortoise-and-hare:
run one pointer at single speed and another at double speed; they are guaranteed to meet inside the
cycle, and you keep only two group elements in memory at a time.

\begin{figure}[H]
\centering
\begin{tikzpicture}[
    pt/.style={circle, draw=structurepurple, fill=structurepurple!12, inner sep=1.3pt},
    walk/.style={-{Stealth[length=2mm]}, structurepurple, thick},
    font=\small
  ]
  % tail
  \node[pt, label=left:{$x_0$}] (t0) at (0,0) {};
  \node[pt] (t1) at (1.1,0) {};
  \node[pt] (t2) at (2.2,0) {};
  % loop (a hexagon-ish cycle)
  \node[pt] (c0) at (3.4,0) {};
  \node[pt] (c1) at (4.4,0.7) {};
  \node[pt] (c2) at (5.2,0) {};
  \node[pt] (c3) at (4.8,-0.9) {};
  \node[pt] (c4) at (3.7,-0.9) {};
  \draw[walk] (t0) -- (t1);
  \draw[walk] (t1) -- (t2);
  \draw[walk] (t2) -- (c0);
  \draw[walk] (c0) -- (c1);
  \draw[walk] (c1) -- (c2);
  \draw[walk] (c2) -- (c3);
  \draw[walk] (c3) -- (c4);
  \draw[walk] (c4) -- (c0);
  \node[font=\scriptsize, structurepurple] at (1.1,0.4) {tail};
  \node[font=\scriptsize, structurepurple] at (4.45,-0.05) {cycle};
  \node[trapred, font=\scriptsize, align=center] at (3.4,0.75) {collision\\$x_i = x_j$};
  \draw[trapred, -{Stealth[length=1.6mm]}, thick] (3.4,0.55) -- (c0);
\end{tikzpicture}
\caption{The $\rho$ shape of Pollard's walk: a tail $x_0,x_1,\dots$ feeds into a cycle. The first
point where the walk re-enters ground it has already visited is the collision $x_i = x_j$ that pins
down the discrete log. Floyd's two pointers meet somewhere inside the loop.}
\label{ch20:fig-rho}
\end{figure}

\begin{whybox}
\textbf{Why $\sqrt{n}$ and not less, generically.} Both BSGS and rho are special cases of a general
truth: in a black-box group where the only thing you can do is multiply and test equality, any
algorithm needs $\Omega(\sqrt{n})$ operations to find a discrete log (Shoup, 1997). The intuition is
the birthday bound running in reverse -- you cannot learn the hidden exponent until two of your probe
elements collide, and random collisions in a set of size $n$ take $\sqrt{n}$ probes. So $\sqrt{n}$ is
not a failure of cleverness; it is a wall for any attack that ignores the internal structure of the
group. The subexponential attacks on $\Zns{p}$ get under the wall precisely by exploiting that
internal structure (index calculus), which is why elliptic curves -- where no such structure is
known -- keep the wall intact and let you use much smaller keys.
\end{whybox}

\begin{example}[Extracting $x$ from a rho collision]
\label{ch20:ex-rho}
The bookkeeping part of rho is what is worth doing by hand, and it goes cleanest in a prime-order
group where every nonzero difference is invertible. Work in a subgroup of prime order
$q = 11$, generator $g$, target $h = g^{x}$. The walk reports a collision between two points carrying
the labels $x_i = g^{8} h^{2}$ and $x_j = g^{3} h^{9}$. Since these are the same group element,
\[
  g^{8} h^{2} = g^{3} h^{9}
  \;\Longrightarrow\;
  g^{8-3} = h^{9-2}
  \;\Longrightarrow\;
  g^{5} = g^{7x},
\]
so $5 \equiv 7x \pmod{11}$. Invert $7$: since $7\cdot 8 = 56 \equiv 1 \pmod{11}$, we have
$7^{-1}\equiv 8$, and $x \equiv 5\cdot 8 = 40 \equiv 7 \pmod{11}$. The walk never needed a table --
two collided labels and one modular inverse pinned the log down. (If $q$ were composite and the
relevant difference shared a factor with it, the inverse would not exist and we would restart the
walk; this is exactly why prime order is so convenient.)
\end{example}

\begin{implementbox}
A compact Pollard rho for discrete logs in $\Zns{p}$, with the exponent bookkeeping spelled out.
\textbf{Learning only -- this has no constant-time discipline, no validation, and is meant to make
the algorithm legible, never to attack or protect anything real.}
\begin{lstlisting}[language=Python]
from math import gcd

def pollard_rho_dlog(g, h, p):
    n = p - 1                      # order of the group Z_p^*
    def step(x, a, b):
        # partition the group into 3 zones by x mod 3
        if   x % 3 == 0:           # zone 0: multiply by g
            return (x*g) % p, (a+1) % n, b
        elif x % 3 == 1:           # zone 1: multiply by h
            return (x*h) % p, a, (b+1) % n
        else:                      # zone 2: square
            return (x*x) % p, (2*a) % n, (2*b) % n
    # tortoise (one step) and hare (two steps)
    x, a, b = 1, 0, 0
    X, A, B = 1, 0, 0
    for _ in range(n):
        x, a, b = step(x, a, b)
        X, A, B = step(*step(X, A, B))
        if x == X:                 # collision found
            r = (b - B) % n        # need to solve  r * dlog = (A - a)  (mod n)
            if r == 0:
                return None        # rare degenerate case; restart with new walk
            t = (A - a) % n
            d = gcd(r, n)          # n is composite here, so d may exceed 1
            if t % d != 0:
                return None        # no solution this run; restart with new walk
            # d candidate residues mod n; exponentiate to pick the right one
            base = (t // d) * pow(r // d, -1, n // d) % (n // d)
            for k in range(d):
                cand = (base + k * (n // d)) % n
                if pow(g, cand, p) == h % p:
                    return cand
            return None
    return None

# Demo (toy parameters): dlog of 17 base 2 mod 29 prints 21.
print(pollard_rho_dlog(2, 17, 29))
\end{lstlisting}
The single modular inverse of the prime-order Example~\ref{ch20:ex-rho} becomes a short \texttt{gcd}
dance here because $n = p-1$ is even (hence composite): the collision difference $r$ can share a
factor with $n$, leaving a handful of candidate logs that one exponentiation each resolves. In a
prime-order subgroup that whole branch collapses to a single inverse -- one more reason cryptographic
deployments live in prime-order subgroups. The cost is $O(\sqrt{n})$ multiplications and a few
integers of memory. For $n$ around $2^{256}$ that is $2^{128}$ multiplications -- the same wall as
BSGS, but now with no storage problem, which is why rho (and its parallel ``van Oorschot-Wiener''
variant) is the practical generic attack of record, and the one used to set elliptic-curve key sizes.
\end{implementbox}

\section{From hard problem to encryption: ElGamal}
\label{ch20:elgamal}

We now have a problem we believe is hard. ElGamal turns it into public-key encryption. The story is
the Diffie-Hellman story with one extra move.

\begin{intuitionbox}
Recall Diffie-Hellman. Alice and Bob agree on $g$ in a public group. Alice picks secret $a$ and
publishes $A = g^a$; Bob picks secret $b$ and publishes $B = g^b$. Each computes the shared secret
$g^{ab} = A^b = B^a$. Eve sees $g, g^a, g^b$ and -- believing CDH -- cannot get $g^{ab}$.

ElGamal's insight: Bob does not need to be online. Alice's published $A = g^a$ \emph{is} her public
key, permanently. Whenever Bob wants to send Alice a secret message $M$, he plays Alice's half of a
\emph{fresh} Diffie-Hellman: he picks a one-time random $y$, forms his half $c_1 = g^y$, computes the
shared secret $s = A^y = g^{ay}$ on the spot, and uses $s$ as a one-time mask to hide $M$, sending
$c_2 = M \cdot s$. Alice, holding $a$, reconstructs the same secret from Bob's half:
$s = c_1^{a} = g^{ya}$, divides it out, and reads $M$. ElGamal is Diffie-Hellman with the second
party replaced by a fresh random exponent for every message.
\end{intuitionbox}

\subsection{The scheme}

\begin{definition}[ElGamal encryption]
\label{ch20:def-elgamal}
Fix a cyclic group $\G = \langle g \rangle$ of prime order $q$ in which DDH is believed hard, and in
which messages are encoded as group elements.
\begin{itemize}
  \item \textbf{Key generation $\Gen$:} choose $a \getsr \{1,\dots,q-1\}$, set $A = g^{a}$. Public
    key $pk = (\G, g, A)$; secret key $sk = a$.
  \item \textbf{Encryption $\Enc_{pk}(M)$:} choose a fresh $y \getsr \{1,\dots,q-1\}$, output the
    ciphertext
    \[
      (c_1, c_2) \;=\; \big(\,g^{y},\;\; M \cdot A^{y}\,\big).
    \]
  \item \textbf{Decryption $\Dec_{sk}(c_1, c_2)$:} compute $s = c_1^{a}$ and recover
    $M = c_2 \cdot s^{-1}$.
\end{itemize}
\end{definition}

\begin{whybox}
\textbf{Correctness.} The shared mask computed by each side agrees:
\[
  c_1^{a} \;=\; (g^{y})^{a} \;=\; g^{ay} \;=\; (g^{a})^{y} \;=\; A^{y}.
\]
So $\Dec$ computes $c_2 \cdot (c_1^a)^{-1} = (M \cdot A^y)\cdot (A^y)^{-1} = M$. The blinding factor
$A^y$ is exactly the Diffie-Hellman secret between the long-term key $A=g^a$ and the one-time key
$c_1 = g^y$; Bob raises $A$ to $y$, Alice raises $c_1$ to $a$, and commutativity of exponents makes
them meet.
\end{whybox}

\begin{example}[ElGamal by hand, $p = 23$]
\label{ch20:ex-elgamal}
We use $\Zns{23}$ for arithmetic that fits in your head. Its order is $22$, which is not prime, so
this toy is not a secure group (more on that below), but the mechanics are identical. Take $g = 5$;
one checks $5$ is a primitive root mod $23$.

\emph{Keygen.} Alice picks $a = 6$. Her public key is $A = 5^{6} \bmod 23$. Compute:
$5^2 = 25 \equiv 2$, $5^3 = 10$, $5^6 = (5^3)^2 = 100 \equiv 8$. So $A = 8$.

\emph{Encrypt $M = 10$.} Bob picks $y = 3$. Then
\[
  c_1 = g^{y} = 5^{3} = 10,\qquad
  c_2 = M \cdot A^{y} = 10 \cdot 8^{3} \bmod 23.
\]
Now $8^2 = 64 \equiv 18$, $8^3 = 8\cdot 18 = 144 \equiv 144 - 138 = 6$. So $c_2 = 10 \cdot 6 = 60
\equiv 60 - 46 = 14$. Ciphertext $(c_1, c_2) = (10, 14)$.

\emph{Decrypt.} Alice computes $s = c_1^{a} = 10^{6} \bmod 23$. Step it: $10^2 = 100 \equiv 8$,
$10^3 = 10\cdot 8 = 80 \equiv 80-69 = 11$, $10^6 = (10^3)^2 = 121 \equiv 121 - 115 = 6$. So $s = 6$,
matching $A^y = 6$ as it must. Invert: $6^{-1} \bmod 23$. Since $6 \cdot 4 = 24 \equiv 1$, we get
$6^{-1} = 4$. Finally $M = c_2 \cdot s^{-1} = 14 \cdot 4 = 56 \equiv 56 - 46 = 10$. Recovered.
Figure~\ref{ch20:fig-elgamal} shows the same exchange schematically.
\end{example}

\begin{figure}[H]
\centering
\begin{tikzpicture}[
    box/.style={draw, rounded corners, align=center, minimum height=8mm, inner sep=4pt, font=\small},
    alice/.style={box, draw=aliceblue, fill=aliceblue!8},
    bob/.style={box, draw=bobgreen, fill=bobgreen!8},
    eve/.style={box, draw=evered, fill=evered!8},
    msg/.style={-{Stealth[length=2.5mm]}, thick}
  ]
  \node[alice] (apk) {Alice publishes\\$A = g^{a}$ (long-term)};
  \node[bob, right=42mm of apk] (benc) {Bob picks fresh $y$\\$c_1 = g^{y}$\\$c_2 = M\cdot A^{y}$};
  \node[alice, below=16mm of apk] (adec) {Alice computes\\$s = c_1^{a} = g^{ay}$\\$M = c_2 / s$};
  \node[eve, below=16mm of benc] (eve) {Eve sees $g, A, c_1, c_2$.\\Needs $g^{ay}$ from $g^a, g^y$:\\that is CDH.};

  \draw[msg, aliceblue] (apk) -- node[above, font=\scriptsize]{public key} (benc);
  \draw[msg, bobgreen] (benc) -- node[right, font=\scriptsize, align=left]{ciphertext\\$(c_1,c_2)$} (adec);
  \draw[msg, dashed, evered] (benc) -- (eve);
\end{tikzpicture}
\caption{ElGamal as a one-sided Diffie-Hellman. Alice's long-term public key supplies one half of the
exchange; Bob's fresh $y$ supplies the other, and the shared secret $g^{ay}$ masks the message. Eve,
seeing only $g^a$ and $g^y$, faces CDH.}
\label{ch20:fig-elgamal}
\end{figure}

\subsection{Why it is secure}

\begin{focusbox}
\textbf{Threat model.} The attacker is a probabilistic polynomial-time eavesdropper who sees the
public key and a ciphertext. The goal we promise to deny her is \emph{indistinguishability under
chosen-plaintext attack} (IND-CPA): she picks two messages $M_0, M_1$, we encrypt one of them at
random, and she cannot guess which better than a coin flip. We do \emph{not} promise IND-CCA security
here -- ElGamal as stated is malleable, as the homomorphic property below makes vivid, and a chosen-
ciphertext attacker breaks it. Real deployments either add a MAC/KEM-DEM wrapper or use a CCA-secure
variant.
\end{focusbox}

\begin{definition}[Decisional Diffie-Hellman]
\label{ch20:def-ddh}
DDH is hard in $\G = \langle g\rangle$ of order $q$ if no $\PPT$ distinguisher can tell apart the two
distributions, over random $a,b,c \in \{0,\dots,q-1\}$,
\[
  \big(g^{a},\, g^{b},\, g^{ab}\big)
  \qquad\text{vs.}\qquad
  \big(g^{a},\, g^{b},\, g^{c}\big),
\]
with non-negligible advantage. In words: a real Diffie-Hellman triple is indistinguishable from one
where the third coordinate is a random group element.
\end{definition}

\begin{theorem}[ElGamal is IND-CPA under DDH]
\label{ch20:thm-ind-cpa}
If DDH is hard in $\G$, then ElGamal over $\G$ is IND-CPA secure.
\end{theorem}

\begin{whybox}
\textbf{Proof sketch (the heart of the argument).} A ciphertext is $(g^y,\ M\cdot A^y)$ with
$A = g^a$. Rewrite it as $(g^y,\ M\cdot g^{ay})$. The second coordinate hides $M$ behind the mask
$g^{ay}$, and the attacker's view is the triple $(g^a, g^y, g^{ay})$ \emph{plus} the message folded
into the third slot.

Now run a hybrid. Replace the genuine mask $g^{ay}$ with a \emph{uniformly random} group element
$g^{c}$. By DDH, no $\PPT$ attacker notices the swap -- the two worlds are computationally
indistinguishable. But in the swapped world the ciphertext is $(g^y,\ M\cdot g^{c})$ with $g^c$
uniform and independent of $M$. A uniform group element times $M$ is again uniform and carries
\emph{zero} information about $M$ -- it is a one-time pad in the group. So in the swapped world the
attacker's advantage is exactly zero, and DDH bounds how much advantage she could have had in the
real world:
\[
  \Adv^{\text{ind-cpa}}_{\text{ElGamal}}(\mathcal{A}) \;\le\; \Adv^{\text{ddh}}_{\G}(\mathcal{B}).
\]
Encryption is secure precisely because the Diffie-Hellman secret $g^{ay}$ looks random to anyone
without $a$ or $y$, and a random group element makes a perfect one-time mask. That is the whole
scheme in one sentence.
\end{whybox}

\begin{trapbox}
\textbf{DDH must really hold -- and in $\Zns{p}$ it often does not!} DDH is \emph{false} in the full
group $\Zns{p}$, because the Legendre symbol (quadratic residue test) leaks one bit of $g^{ab}$ that
an attacker can predict from $g^a$ and $g^b$. Concretely, whether $g^{ab}$ is a quadratic residue is
determined by whether $a$ and $b$ are even, which the QR-ness of $g^a, g^b$ reveals. A distinguisher
that just compares Legendre symbols wins. \textbf{Fix:} work in the prime-order subgroup of quadratic
residues -- the order-$q$ subgroup of $\Zns{p}$ where $p = 2q+1$ -- in which every element is a QR and
the leak vanishes. This is why ``ElGamal over $\Zns{p}$'' done right means ``ElGamal in the
prime-order subgroup,'' not the whole group.
\end{trapbox}

\subsection{The homomorphic property}
\label{ch20:homomorphic}

ElGamal has a structural feature that is sometimes a bug and sometimes the entire point:
multiplying two ciphertexts encrypts the product of the plaintexts. Take two encryptions of $M$ and
$M'$ under the same key, with independent randomness $y, y'$:
\[
  (c_1, c_2) = (g^{y},\ M A^{y}),
  \qquad
  (c_1', c_2') = (g^{y'},\ M' A^{y'}).
\]
Multiply componentwise:
\[
  (c_1 c_1',\ c_2 c_2')
  \;=\;
  \big(g^{y+y'},\ (M M')\,A^{y+y'}\big),
\]
which is a valid ElGamal encryption of $M M'$ under the same key, with randomness $y + y'$. The
scheme is \emph{multiplicatively homomorphic}.

\begin{intuitionbox}
You can compute on the data without ever decrypting it. A tallying server can multiply together
ciphertexts to combine encrypted values, and only the keyholder, at the end, learns the combined
result -- never the individual inputs. This is the seed of homomorphic encryption (its own chapter
later) and of cryptographic voting. By tweaking the scheme to encrypt $g^{M}$ instead of $M$
(``exponential ElGamal''), multiplication of ciphertexts becomes \emph{addition} of the messages in
the exponent, which is exactly what you want for summing votes -- at the price of having to solve a
small discrete log at the end to read the tally off $g^{\sum M_i}$.
\end{intuitionbox}

\begin{attackbox}
\textbf{Malleability: the flip side.} The same property breaks chosen-ciphertext security. Suppose
Eve intercepts a ciphertext $(c_1, c_2)$ for an unknown price $M$ and wants to change it to encrypt
$2M$. She does not need the key: she sends $(c_1,\ 2 c_2)$, and the recipient decrypts
$2 c_2 \cdot s^{-1} = 2M$. She mauled the plaintext blind. Worse, against a decryption oracle she can
recover $M$ entirely: ask the oracle to decrypt $(c_1,\ r\cdot c_2)$ for a known random $r$; it
returns $r M$; divide by $r$ to get $M$. This is why textbook ElGamal is never used raw for
confidentiality. Production systems wrap it: hash the Diffie-Hellman secret into a symmetric key and
encrypt-then-MAC the payload (the KEM/DEM pattern of ECIES), which removes malleability by binding the
ciphertext together under authentication.
\end{attackbox}

\section{Choosing safe groups and parameters}
\label{ch20:params}

Getting ElGamal (or Diffie-Hellman) right is mostly about the group. Three rules.

\textbf{1. Use a prime-order subgroup.} Work in a subgroup of prime order $q$, not in $\Zns{p}$
whole. The standard construction: pick a \emph{safe prime} $p = 2q + 1$ with $q$ also prime. The
quadratic residues mod $p$ then form a subgroup of order $q$, and you take $g$ to be a generator of
\emph{that} subgroup. Prime order is the gift that keeps giving: every non-identity element is a
generator, there are no small subgroups to fall into, the Pohlig-Hellman attack (which would split
the DLP across the prime factors of the order) has nothing to bite on, and DDH's Legendre-symbol leak
disappears.

\begin{focusbox}
\textbf{Why prime order kills small-subgroup attacks.} If the order $n$ factors as $n = q_1 q_2
\cdots$, Pohlig-Hellman solves the DLP separately modulo each $q_i$ -- in time $\sqrt{q_{\max}}$,
governed by the \emph{largest} factor -- and reassembles by CRT. A composite order with all-small
factors is catastrophic: order $2^{40} \cdot (\text{small primes})$ would be breakable on a laptop
regardless of how big $p$ is. Prime order $q$ means there is exactly one prime factor and it is the
whole thing, so $\sqrt{q}$ is the true cost. This is the discrete-log analogue of ``a chain is as
strong as its weakest link.''
\end{focusbox}

\textbf{2. Size for the best attack, on both fronts.} Two clocks tick. The generic attack (rho)
costs $\sqrt{q}$, so $q$ must be at least $256$ bits to reach the $2^{128}$ generic-security level.
The index-calculus attack costs $L_p[1/3,\cdot]$ in the size of $p$, so $p$ must be far larger --
about $3072$ bits for the same $128$-bit target, and the older $1024$-bit primes are now considered
broken by well-resourced adversaries. You size $p$ against index calculus and $q$ against rho;
the safe-prime construction $p = 2q+1$ couples them, which is why it forces a large $q$ along with
a large $p$.

\textbf{3. Fresh, uniform randomness for $y$ every time.} The per-message exponent $y$ must be
freshly sampled and never reused.

\begin{attackbox}
\textbf{Nonce reuse murders ElGamal.} Suppose Bob lazily reuses the same $y$ for two messages
$M_1, M_2$. The ciphertexts are $(g^y, M_1 A^y)$ and $(g^y, M_2 A^y)$ -- same $c_1$, so an
eavesdropper instantly spots the reuse, and the ratio of second components is
\[
  \frac{M_1 A^y}{M_2 A^y} \;=\; \frac{M_1}{M_2},
\]
leaking the ratio of the two plaintexts with the key entirely irrelevant. If she ever learns $M_1$
(a known header, a guessed field), she gets $M_2$ for free. This is the same lesson as IV reuse in
stream ciphers and nonce reuse in GCM, and it is a recurring theme: the randomness in a randomized
scheme is not optional flavor, it is load-bearing structure. The signature cousin of ElGamal (DSA) is
even less forgiving -- reusing $y$ there leaks the \emph{secret key}, the bug that infamously
unlocked the Sony PlayStation 3.
\end{attackbox}

\begin{connectionbox}
The cost of the \emph{forward} direction -- computing $g^y$, $A^y$, $c_1^a$ -- is square-and-multiply
modular exponentiation, $O(\log q)$ multiplications, each on numbers of size $\log p$. That is the
algorithm analyzed in the modular-exponentiation section of the algorithms companion. The asymmetry
that makes ElGamal possible is exactly the gap between this $O(\log q)$ forward cost and the
$O(\sqrt{q})$ (or subexponential, via index calculus) backward cost. Public-key cryptography lives in
that gap.
\end{connectionbox}

\begin{applicationbox}
\textbf{Where this actually runs.} Pure mod-$p$ ElGamal is rare in modern transport security -- it is
slow and the ciphertext doubles the message size -- but its DNA is everywhere:
\begin{itemize}
  \item \textbf{GnuPG / OpenPGP} has long offered ElGamal as an encryption key type (paired with DSA
    signatures), encrypting a random session key that then encrypts the actual email or file under a
    fast symmetric cipher -- the hybrid pattern.
  \item \textbf{ECIES}, the Elliptic Curve Integrated Encryption Scheme, is ElGamal moved onto an
    elliptic-curve group and hardened: derive a symmetric key from the Diffie-Hellman secret with a
    KDF, then encrypt-then-MAC. It is the standard way to do public-key encryption to a recipient on
    a curve, used in messaging and in cryptocurrency wallets.
  \item \textbf{TLS} does not encrypt with ElGamal, but its ephemeral key exchange (DHE and ECDHE) is
    the Diffie-Hellman half of this story, run fresh per session for forward secrecy -- the subject
    of the key-exchange chapter.
  \item \textbf{Cryptographic voting} (Helios and relatives) uses \emph{exponential} ElGamal exactly
    for its additive homomorphism: encrypted ballots are multiplied to tally without decrypting
    individual votes.
\end{itemize}
\end{applicationbox}

\section{Relationship to Diffie-Hellman}
\label{ch20:dh-relationship}

It is worth stating plainly how the three problems and the two schemes line up, because they are
constantly conflated.

\begin{itemize}
  \item \textbf{Diffie-Hellman key exchange} is an interactive protocol: both parties contribute a
    fresh ephemeral exponent, they jointly compute $g^{ab}$, and they get a shared symmetric key.
    Nothing is ``encrypted''; a key is \emph{agreed}. Its security rests on CDH (to hide the key) and
    on DDH (to argue the key looks uniform).
  \item \textbf{ElGamal encryption} is non-interactive: one party (the recipient) has a long-lived
    key $g^a$; the sender runs a one-sided Diffie-Hellman with a fresh $y$ and folds the message into
    the resulting secret. Security rests on DDH (Theorem~\ref{ch20:thm-ind-cpa}).
\end{itemize}

So ElGamal \emph{is} Diffie-Hellman, frozen into a public-key encryption scheme by making one side's
contribution permanent and the other side's per-message. The shared secret $g^{ay}$ plays in ElGamal
exactly the role $g^{ab}$ plays in DH. And the same hardness assumptions underwrite both: break the
DLP and you break everything; break only CDH and you can read DH-keyed traffic and forge ElGamal
masks; break only DDH and you cannot necessarily recover keys but you \emph{can} distinguish
ciphertexts, which is enough to demolish the IND-CPA guarantee.

\begin{focusbox}
\textbf{One picture to keep.} Forward ($g, x \mapsto g^x$) is cheap; backward ($g, g^x \mapsto x$) is
the DLP and is hard. Diffie-Hellman builds a shared secret in that gap; ElGamal freezes one half of
the exchange into a public key and masks a message with the secret. The whole edifice -- and its
elliptic-curve descendants in the next chapter -- is the discrete logarithm problem wearing different
costumes.
\end{focusbox}

\section{Exercises}
\label{ch20:exercises}

\begin{exercise}[Hand-cranking discrete logs]
\label{ch20:exc-table}
Work in $\Zns{13}$ with $g = 2$. (a) Verify $2$ is a primitive root by listing all powers
$2^0,\dots,2^{11} \bmod 13$. (b) Read off $\dlog_2 6$, $\dlog_2 7$, and $\dlog_2 11$ from your table.
(c) Confirm the identity $\dlog_2(6\cdot 7) \equiv \dlog_2 6 + \dlog_2 7 \pmod{12}$, and explain in
one line why this additivity must hold.
\end{exercise}

\begin{exercise}[Baby-step giant-step]
\label{ch20:exc-bsgs}
Let $p = 41$, $g = 6$ (a primitive root mod $41$). Use baby-step giant-step with $m = \ceil{\sqrt{40}}
= 7$ to compute $\dlog_6 29$. Show the baby-step table, the giant-step value $g^{-m} \bmod 41$, and
the iteration that produces the hit. Verify your answer by direct exponentiation.
\end{exercise}

\begin{exercise}[Pollard rho extraction]
\label{ch20:exc-rho}
In a group of prime order $q = 101$ with generator $g$, a rho walk reports a collision between
$x_i = g^{12} h^{7}$ and $x_j = g^{83} h^{4}$, where $h = g^{x}$ is the unknown target. Solve the
linear congruence for $x \bmod 101$. (You will need $3^{-1} \bmod 101$; find it.) Why was it
essential that the order be prime for this step to go through cleanly?
\end{exercise}

\begin{exercise}[ElGamal round trip]
\label{ch20:exc-elgamal}
Use $\Zns{23}$, $g = 5$, secret key $a = 9$. (a) Compute Alice's public key $A$. (b) Encrypt
$M = 15$ using one-time randomness $y = 4$, giving $(c_1, c_2)$. (c) Decrypt your ciphertext as Alice
and confirm you recover $15$. Show every modular reduction.
\end{exercise}

\begin{exercise}[Homomorphic combination]
\label{ch20:exc-homo}
Continuing in $\Zns{23}$ with $A$ from the previous exercise: encrypt $M_1 = 3$ with $y_1 = 2$ and
$M_2 = 4$ with $y_2 = 5$. Multiply the two ciphertexts componentwise, decrypt the product ciphertext,
and verify it yields $M_1 M_2 = 12 \bmod 23$. Then state what randomness the product ciphertext
``uses,'' and confirm it is $y_1 + y_2$.
\end{exercise}

\begin{exercise}[Nonce reuse]
\label{ch20:exc-nonce}
Bob encrypts two messages to Alice under the same public key and accidentally reuses the same $y$,
producing $(c_1, c_2)$ and $(c_1, c_2')$. (a) How can an eavesdropper detect the reuse from the
ciphertexts alone? (b) Show that the eavesdropper learns $M / M'$ without the secret key. (c) If a
known header tells her $M = 1000$ and she computes $M/M' = 8$, what is $M'$? (d) Contrast with DSA,
where reusing the per-signature nonce leaks something far worse -- name what.
\end{exercise}

\begin{exercise}[The DDH leak in the full group]
\label{ch20:exc-legendre}
Let $p$ be an odd prime and $g$ a primitive root, so $g$ is a quadratic non-residue. (a) Show that
$g^{x}$ is a quadratic residue if and only if $x$ is even. (b) Given $g^a$ and $g^b$, explain how to
predict the QR-ness of $g^{ab}$, and use this to build a DDH distinguisher with advantage $1/4$ in
the full group $\Zns{p}$. (c) Explain why restricting to the order-$q$ subgroup of quadratic residues
(with $p = 2q+1$) shuts this distinguisher down.
\end{exercise}

\begin{exercise}[Pohlig-Hellman, why prime order matters]
\label{ch20:exc-pohlig}
Suppose a group has \emph{smooth} order $n = 2 \cdot 3 \cdot 5 \cdot 7 = 210$. Sketch how
Pohlig-Hellman reduces one discrete log in this group to four small discrete logs (mod $2,3,5,7$),
each solvable by inspection, recombined by the Chinese Remainder Theorem. Estimate the total work and
compare it to $\sqrt{210}$. What does this say about ever using a group whose order has only small
prime factors?
\end{exercise}

\begin{exercise}[Malleability into a CCA break]
\label{ch20:exc-cca}
You are given an ElGamal decryption oracle (it decrypts anything except the exact challenge
ciphertext $(c_1, c_2)$) and the challenge ciphertext itself, encrypting unknown $M$. Design a single
oracle query that recovers $M$, using the multiplicative homomorphism. Why does encrypt-then-MAC over
the ciphertext defeat your attack?
\end{exercise}

\begin{exercise}[Sizing a real parameter set]
\label{ch20:exc-sizing}
You want $128$-bit security for ElGamal over $\Zns{p}$ with a safe prime $p = 2q + 1$. (a) Using the
generic $\sqrt{q}$ bound, what is the minimum bit-length of $q$? (b) Given that index calculus on $p$
costs roughly $L_p[1/3]$, mainstream guidance puts $p$ at about $3072$ bits for this level -- explain
why $p$ must be \emph{so much} larger than $q$. (c) An elliptic curve achieves the same $128$-bit
security with a $256$-bit field. In one or two sentences, what structural difference lets the curve
get away with the smaller parameter?
\end{exercise}

\begin{exercise}[Prove the homomorphism]
\label{ch20:exc-prove-homo}
Prove from the definitions that ElGamal is multiplicatively homomorphic: if $(c_1, c_2)$ encrypts $M$
with randomness $y$ and $(c_1', c_2')$ encrypts $M'$ with randomness $y'$, then
$(c_1 c_1',\, c_2 c_2')$ is a valid encryption of $M M'$, and identify its randomness. Then prove the
``exponential ElGamal'' claim: if instead you encrypt $g^{M}$ and $g^{M'}$, the product ciphertext
decrypts to $g^{M + M'}$, so reading off $M + M'$ requires one discrete log.
\end{exercise}
